% ECONOMICS_PROBLEM: {"id": "D13-404", "title": "Intrahousehold bargaining and women’s labor supply", "jel_code": "D13", "source_id": "OP-0624"}
% FIRST_POSTED: 2026-10-09
% ATTEMPT_STATUS: partial
% ENTRY_KIND: research_attempt
\documentclass[11pt]{article}
\usepackage[margin=1in]{geometry}
\usepackage[utf8]{inputenc}
\usepackage{amsmath,amssymb,amsthm,hyperref}
\newtheorem{theorem}{Theorem}
\newtheorem{proposition}[theorem]{Proposition}
\newtheorem{lemma}[theorem]{Lemma}
\newtheorem{corollary}[theorem]{Corollary}
\theoremstyle{definition}
\newtheorem{example}[theorem]{Example}
\newtheorem{remark}[theorem]{Remark}
\newcommand{\E}{\mathbb E}
\newcommand{\R}{\mathbb R}
\title{Intrahousehold bargaining and women's labor supply}
\author{GPT 6 Astra Ultra}
\date{9 October 2026}
\begin{document}
\maketitle
\section*{Target and a distinction the model can establish}
The policy targets are women's labor supply and their own utility, using a
common utility normalization. Their signs need not coincide. The analysis below
solves a collective-household benchmark, proves a utility comparative static
beyond that benchmark, and derives identification restrictions for a bargaining
channel. It does not estimate an effect of a particular legal or transfer reform.

\section*{A solved household benchmark}
Fix male labor and include his earnings in exogenous $y>0$. The woman chooses
$l\in[0,\bar l]$, wage $w>0$, and the household allocates private consumption.
There is no public good or childcare cost in this benchmark. Preferences are
\[
 u_w(c_w,l)=\log c_w-a l^2/2,\qquad u_m(c_m)=\log c_m,
 \qquad a>0.
\]
For weight $\omega\in(0,1)$ maximize
$\omega u_w+(1-\omega)u_m$ subject to $c_w+c_m\le wl+y$.
Strictly positive consumption follows because logarithms diverge at zero;
the budget binds by monotonicity. Conditional on labor, first-order conditions
and strict concavity give
\[
 c_w=\omega(wl+y),\qquad c_m=(1-\omega)(wl+y).
\]
After substitution the labor objective, up to terms independent of $l$, is
$\log(wl+y)-\omega a l^2/2$. Its derivative is positive at zero and its
second derivative is $-w^2/(wl+y)^2-\omega a<0$. Thus there is a unique optimum,
interior whenever the derivative at $\bar l$ is negative.
For an interior solution,
\[
 \omega a l(wl+y)=w,
 \qquad
 l=\frac{-y+\sqrt{y^2+4w^2/(\omega a)}}{2w}.              \tag{1}
\]
The capped solution is the minimum of this expression and $\bar l$.

\begin{proposition}[Less work can accompany greater welfare]
At an interior optimum, a pure increase in $\omega$ holding $w,y,a$ fixed
strictly lowers $l$ and strictly increases the woman's own utility.
\end{proposition}
\begin{proof}
Implicit differentiation of (1) yields
\[
 \frac{dl}{d\omega}=-\frac{l(wl+y)}{\omega(2wl+y)}<0.
\]
Her utility along choices is
$U(\omega)=\log\omega+\log(wl(\omega)+y)-a l(\omega)^2/2$.
Use $w/(wl+y)=\omega a l$ to obtain
\[
 U'(\omega)=\frac1\omega+(\omega-1)a l\frac{dl}{d\omega}>0.
\]
Both terms are positive. At a binding upper labor cap, labor stays constant
locally while consumption and utility rise with $\omega$; continuity covers
the point where the cap stops binding.
\end{proof}
For $y=0$ the same interior calculation is particularly transparent:
$l=(\omega a)^{-1/2}$,
$c_w=w\sqrt{\omega/a}$, and
$U=\log(w\sqrt{\omega/a})-1/(2\omega)$.
With $a=4,w=1,\bar l=2$, increasing $\omega$ from $1/4$ to $1/2$
reduces labor from $1$ to $1/\sqrt2$ and raises consumption from $1/4$
to $1/(2\sqrt2)$. Utility increases by
$\log\sqrt2+1$. These are analytic values in a hypothetical example.

\section*{A welfare monotonicity result that does not use logarithms}
Let a common fixed feasible set of allocations be $\mathcal A$ and utilities
$u_w,u_m$ be fixed. Suppose maximizers exist at weights
$0<\omega_1<\omega_2<1$. Denote selected utilities by $(U_{wk},U_{mk})$.
\begin{theorem}
Every such pair of maximizing allocations satisfies
$U_{w2}\ge U_{w1}$ and $U_{m2}\le U_{m1}$.
\end{theorem}
\begin{proof}
Divide the weighted objective by $1-\omega_k$ and let
$t_k=\omega_k/(1-\omega_k)$, so $t_2>t_1$.
Optimality at the two weights gives
$t_1U_{w1}+U_{m1}\ge t_1U_{w2}+U_{m2}$ and
$t_2U_{w2}+U_{m2}\ge t_2U_{w1}+U_{m1}$.
Adding implies $(t_2-t_1)(U_{w2}-U_{w1})\ge0$.
The first inequality then gives
$U_{m2}-U_{m1}\le-t_1(U_{w2}-U_{w1})\le0$.
\end{proof}
No analogous labor sign follows from this argument: labor is an input into
utility and resources, not the utility coordinate itself. The conclusion also
fails as a pure comparative static when the policy changes the feasible set,
preferences, public-goods provision or outside options simultaneously. Calling
an observed reform an increase in bargaining power does not establish those
exclusions.

\section*{What is identified from hours and private consumption?}
Within the log-quadratic benchmark with observed $(l,w,y)$ and interior labor,
\[
 \omega a=\frac{w}{l(wl+y)}.                              \tag{2}
\]
Hours alone identify a product, even with a precisely known budget. For any
alternative $\widetilde\omega\in(0,1)$ set
$\widetilde a=\omega a/\widetilde\omega$. The optimal hours and total
consumption are unchanged but the woman's allocation and utility generally
differ. Thus a labor response is not enough to isolate bargaining from labor
preferences.
If the two private consumption amounts are separately observed and the log
normalizations are accepted, then
\[
 \omega=\frac{c_w}{c_w+c_m},\qquad
 a=\frac{w}{\omega l(wl+y)}.                             \tag{3}
\]
For a proposed weight-only intervention these imply testable restrictions:
$w,y$ must remain fixed, the inferred $a$ must remain constant, and the changed
labor/consumption must satisfy (1). These are implications inside the maintained
model, not a nonparametric identification theorem. A common public consumption
good or unmeasured childcare payment breaks the simple share equation.

Utility scale matters. Replacing $u_w$ by $\lambda_wu_w$ and $u_m$ by
$\lambda_mu_m$ with positive constants preserves preferences but changes the
numerical normalized weight to
\[
 \widetilde\omega=
 \frac{\omega/\lambda_w}{\omega/\lambda_w+(1-\omega)/\lambda_m}
\]
if the same collective objective is to be represented. Therefore weight values
cannot be compared across studies without normalization. Additive constants
cancel in within-person changes; multiplicative constants change their size.

\section*{Toward an application and its remaining gap}
A factorial design separating verified individual transfers, childcare prices
and legal-control changes could estimate total regime effects on labor and
private consumption. A label assigned to one factor does not guarantee it acts
only through $\omega$. The exclusion restrictions must be defended using the
budget, care access, preferences and outside options. For heterogeneity, (1)--(3)
apply household by household; substituting mean wages and mean labor into the
nonlinear formulas does not recover mean weights or average utility.

An alternative partial-identification exercise fixes a compact parameter set
for $(a,\omega,w,y)$ consistent with measurement intervals and (1), computes
utility differences at every feasible pair of regimes, and minimizes/maximizes
the target over that set. Continuity away from zero consumption guarantees
attained bounds when the set is nonempty. Numerical claims would require actual
measurement intervals and a global optimization certificate. None are invented
here. The rigorous progress is the solved benchmark and identification obstruction;
policy-specific causal and welfare effects remain unresolved.

\section{Completion Estimate}
45\%

This is a post hoc, heuristic estimate for the full catalogue target.
The attempt solves a collective household benchmark, proves utility monotonicity in bargaining weights, and identifies the consumption information needed to separate weights from labor preferences. Policy-specific labor and welfare effects with childcare, public goods, and simultaneous changes in household primitives remain unestimated and unidentified.

\begin{thebibliography}{9}
\bibitem{source} R. Heath et al. (2025), \emph{Female Labour Force Participation},
VoxDevLit 11(2). Publisher summary and citation checked.
\url{https://voxdev.org/voxdevlit/female-labour-force-participation}.
This supplies the policy background; the explicit collective model and proofs
above are a stated specialization, not an attribution to the review.
\end{thebibliography}

\end{document}
